% !TEX TS-program = pdflatex
% !TEX encoding = UTF-8 Unicode

% This file is a template using the "beamer" package to create slides for a talk or presentation
% - Giving a talk on some subject.
% - The talk is between 15min and 45min long.
% - Style is ornate.

% MODIFIED by Jonathan Kew, 2008-07-06
% The header comments and encoding in this file were modified for inclusion with TeXworks.
% The content is otherwise unchanged from the original distributed with the beamer package.

\documentclass{beamer}
\setbeamercovered{invisible}



% Copyright 2004 by Till Tantau <tantau@users.sourceforge.net>.
%
% In principle, this file can be redistributed and/or modified under
% the terms of the GNU Public License, version 2.
%
% However, this file is supposed to be a template to be modified
% for your own needs. For this reason, if you use this file as a
% template and not specifically distribute it as part of a another
% package/program, I grant the extra permission to freely copy and
% modify this file as you see fit and even to delete this copyright
% notice. 


\mode<presentation>
{
  \usetheme{Warsaw}
  % or ...

  % or whatever (possibly just delete it)
}


\usepackage[english]{babel}
% or whatever

\usepackage[utf8]{inputenc}
\usepackage{qrcode}
\usepackage{hyperref}
\usepackage{ytableau}
\usepackage{longtable}
% or whatever

\usepackage{times}
\usepackage[T1]{fontenc}
% Or whatever. Note that the encoding and the font should match. If T1
% does not look nice, try deleting the line with the fontenc.

%%% MATH RELATED

\title{MATH 180 Strategies of Problem Solving} 
\subtitle
{} % (optional)

\author[W.R. Casper] % (optional, use only with lots of authors)
{W.R. Casper}
% - Use the \inst{?} command only if the authors have different
%   affiliation.

\institute[California State University Fullerton] % (optional, but mostly needed)
{
  Department of Mathematics\\
  California State University Fullerton}
% - Use the \inst command only if there are several affiliations.
% - Keep it simple, no one is interested in your street address.

\subject{Talks}
% This is only inserted into the PDF information catalog. Can be left
% out. 



% If you have a file called "university-logo-filename.xxx", where xxx
% is a graphic format that can be processed by latex or pdflatex,
% resp., then you can add a logo as follows:

% \pgfdeclareimage[height=0.5cm]{university-logo}{university-logo-filename}
% \logo{\pgfuseimage{university-logo}}



% Delete this, if you do not want the table of contents to pop up at
% the beginning of each subsection:
\AtBeginSubsection[]
{
  \begin{frame}<beamer>{Outline}
    \tableofcontents[currentsection,currentsubsection]
  \end{frame}
}


% If you wish to uncover everything in a step-wise fashion, uncomment
% the following command: 

%\beamerdefaultoverlayspecification{<+->}


\begin{document}

\begin{frame}
  \titlepage
\end{frame}

\begin{frame}{Outline}
  \tableofcontents
  % You might wish to add the option [pausesections]
\end{frame}

% Since this a solution template for a generic talk, very little can
% be said about how it should be structured. However, the talk length
% of between 15min and 45min and the theme suggest that you stick to
% the following rules:  

% - Exactly two or three sections (other than the summary).
% - At *most* three subsections per section.
% - Talk about 30s to 2min per frame. So there should be between about
%   15 and 30 frames, all told.

\section{SoPS Lecture 10}
\subsection{Sprouts}

\begin{frame}{Review}
  \begin{block}{Last time}
    \begin{itemize}
\pause
      \item games focused on deductive / spatial reasoning
    \end{itemize}
  \end{block}
\end{frame}

\begin{frame}{Sprouts Game}
  \begin{block}{Rules}
    \begin{itemize}
\pause
      \item draw a handful of \textbf{seeds}, ie. dots, on a piece of paper
\pause
      \item a move consists of drawing a \textbf{sprout}, ie. curve, from one seed to another, and then drawing a new seed somewhere on the sprout
\pause
      \item each seed can have at most three sprouts coming out of it
\pause
      \item sprouts cannot cross other sprouts
\pause
      \item last player to draw an arc wins
    \end{itemize}
  \end{block}
  \begin{block}{Activity (20 minutes)}
    Get into groups of two or three and try playing Sprouts!
  \end{block}
\end{frame}
    
\begin{frame}{Observations:}
  \begin{block}{Question:}
    What kinds of observations can you make about the game Sprouts?
  \end{block}
\pause
  \begin{block}{Question:}
    What kinds of conjectures can you make about the game Sprouts?
  \end{block}
\end{frame}

\begin{frame}{A single seed:}
  \begin{block}{Question}
    What happens if the initial setup has only one seed?
  \end{block}
\end{frame}

\begin{frame}{Two seeds:}
  \begin{block}{Question}
    What about if we start with two seeds, what can happen?
  \end{block}
\end{frame}


\subsection{Graphs}

\begin{frame}{Vocabulary}
\pause
  \begin{block}{Graphs}
    A \textbf{graph} or \textbf{network} is a collection of dots called \textbf{vertices} and arcs between the dots called \textbf{edges}.
  \end{block}
  \begin{block}{Planar graphs}
    A graph is called \textbf{planar} if it can be drawn without the edges crossing each other.
  \end{block}
\end{frame}

\begin{frame}{Game length}
  \begin{block}{Question 1}
    Does every game of Sprouts eventually end?  If we start with $n$ seeds, what's the maximum number of moves?
  \end{block}
  \begin{block}{Question 2}
    When the game has ended, how many regions are enclosed by the final graph?
  \end{block}
\end{frame}

\subsection{Wrapping up}

\begin{frame}{Discussion}
\pause
  \begin{block}{Good question}
    What kind of interesting questions can you think of about Break the Cube?
  \end{block}
\pause
  \begin{block}{Question}
    Suppose we wanted to try to count the number of possible arrangements of the three blocks.  How might we begin to go about it?
  \end{block}
\end{frame}


\begin{frame}{Question of the Day}
\pause
  \begin{block}{Puzzle}
	Suppose there is an enemy submarine at an unknown location on the integer
	number line, and each minute, it moves at an unknown, fixed integer
	velocity.  You may fire one missile per minute somewhere on the number line
	in an attempt to hit the submarine. Is there a strategy to guarantee you
    will eventually hit the submarine?
  \end{block}
\end{frame}

\begin{frame}{Final thoughts}
  \begin{block}{Reflection}
    \begin{itemize}
\pause
      \item What was the static pattern you found that worked the best?
\pause
      \item What dynamic algorithm would you suggest to improve your static algorithm.
    \end{itemize}
  \end{block}
\begin{center}
Submit your responses by \textbf{Tonight!}
\end{center}
\end{frame}


\end{document}


